\section{Propiedades del producto}

Hemos introducido el producto de números naturales mediante sumas repetidas. Esto nos permite hacer algo importante: muchas de sus propiedades pueden comprenderse utilizando únicamente ideas que ya conocemos sobre la adición.

No se trata de memorizar una nueva colección de reglas independientes. En muchos casos podremos desarmar un producto en sumas, utilizar las propiedades de la adición y volver a escribir el resultado como un producto.

\subsection{Clausura}

Como el producto entre números naturales puede construirse mediante sumas de números naturales, su resultado vuelve a ser un número natural.

Por ejemplo,
\[
4\cdot3
=
3+3+3+3
=
12
\]
y
\[
12\in\mathbb{N}
\]
De manera general, si
\[
a,b\in\mathbb{N}
\]
entonces
\[
a\cdot b\in\mathbb{N}
\]
Por eso decimos que los números naturales son \emph{cerrados} respecto del producto.

\subsection{Propiedad conmutativa}

Considere
\[
3\cdot4
\]
Según nuestra interpretación del producto,
\[
3\cdot4
=
4+4+4
=
12.
\]
En cambio,
\[
4\cdot3
=
3+3+3+3
=
12.
\]
En ambos casos obtenemos el mismo número:
\[
3\cdot4=4\cdot3
\]
Este hecho puede comprenderse de una manera sencilla si imaginamos una disposición rectangular de objetos. El producto
\[
3\cdot4
\]
puede representar tres filas de cuatro objetos:
\[
\begin{array}{cccc}
\bullet&\bullet&\bullet&\bullet\\
\bullet&\bullet&\bullet&\bullet\\
\bullet&\bullet&\bullet&\bullet
\end{array}
\]
Pero esa misma disposición también puede observarse como cuatro columnas de tres objetos. Contar por filas o contar por columnas no modifica la cantidad de objetos.
De manera general,
\[
a\cdot b=b\cdot a
\]
Esta propiedad recibe el nombre de \emph{propiedad conmutativa del producto}.
Gracias a ella, aunque inicialmente interpretemos
\[
a\cdot b
\]
como \(a\) repeticiones de \(b\), también podemos interpretarlo como \(b\) repeticiones de \(a\).

\subsection{Elemento neutro}

Considere ahora el producto
\[
1\cdot a
\]
Según la interpretación de suma repetida, esto significa tomar una sola vez el número \(a\). Por lo tanto,
\[
1\cdot a=a
\]
Como el producto es conmutativo,
\[
a\cdot1=1\cdot a=a
\]
Así,
\[
a\cdot1=1\cdot a=a
\]
El número \(1\) recibe el nombre de \emph{elemento neutro del producto}: multiplicar un número por \(1\) no modifica la cantidad representada.

\subsection{Producto por cero}

Veamos ahora qué ocurre con el número \(0\).

La expresión
\[
a\cdot0
\]
representa la suma de \(a\) ceros:
\[
\underbrace{0+0+\cdots+0}_{a\text{ veces}}
\]
Como sumar ceros no modifica la cantidad,
\[
a\cdot0=0
\]
Por la propiedad conmutativa,
\[
0\cdot a=a\cdot0=0
\]
Por lo tanto,
\[
a\cdot0=0\cdot a=0
\]
El cero no es un elemento neutro del producto. Por el contrario, cualquier producto que tenga al \(0\) como factor representa el número \(0\).

\subsection{Propiedad asociativa}

El producto, al igual que la suma, es una operación binaria. Por esta razón, una expresión como
\[
2\cdot3\cdot4
\]
representa aplicaciones sucesivas del producto.
Podríamos agrupar primero
\[
2\cdot3
\]
y escribir
\[
(2\cdot3)\cdot4
\]
Como
\[
2\cdot3=6
\]
obtenemos
\[
6\cdot4=24
\]
Pero también podríamos comenzar por
\[
3\cdot4
\]
y escribir
\[
2\cdot(3\cdot4)
\]
Como
\[
3\cdot4=12
\]
obtenemos
\[
2\cdot12=24
\]
Ambas formas de agrupar representan el mismo número:
\[
(2\cdot3)\cdot4
=
2\cdot(3\cdot4)
\]
Podemos comprender por qué ocurre esto utilizando nuevamente la suma repetida.
La expresión
\[
a\cdot(b\cdot c)
\]
significa tomar \(a\) veces la cantidad \(b\cdot c\):
\[
\underbrace{(b\cdot c)+(b\cdot c)+\cdots+(b\cdot c)}_{a\text{ veces}}
\]
Pero cada \(b\cdot c\) representa, a su vez,
\[
\underbrace{c+c+\cdots+c}_{b\text{ veces}}
\]
Por lo tanto, en total estamos sumando \(c\) una cantidad de \(a\cdot b\) veces. Esto puede escribirse como
\[
(a\cdot b)\cdot c
\]
Así obtenemos, de manera general,
\[
(a\cdot b)\cdot c
=
a\cdot(b\cdot c)
\]
Esta es la \emph{propiedad asociativa del producto}.
Gracias a ella podemos escribir simplemente
\[
a\cdot b\cdot c
\]
sin indicar mediante paréntesis qué producto debe realizarse primero, siempre que en la expresión intervengan únicamente productos.

\subsection{Propiedad distributiva}

Existe además una propiedad especialmente importante porque relaciona las dos operaciones que conocemos hasta ahora: la suma y el producto.

Considere
\[
3\cdot(2+4)
\]
Como el primer factor indica cuántas veces repetimos el segundo,
\[
3\cdot(2+4)
=
(2+4)+(2+4)+(2+4)
\]
Utilizando las propiedades asociativa y conmutativa de la adición podemos reorganizar los sumandos:
\[
(2+2+2)+(4+4+4)
\]
Ahora reconocemos nuevamente dos productos:
\[
2+2+2=3\cdot2
\]
y
\[
4+4+4=3\cdot4
\]
Por lo tanto,
\[
3\cdot(2+4)
=
3\cdot2+3\cdot4
\]
De manera general,
\[
a\cdot(b+c)
=
a\cdot b+a\cdot c
\]
Esta propiedad recibe el nombre de \emph{propiedad distributiva del producto respecto de la suma}.

También podemos encontrar la suma en el primer factor. Por ejemplo,
\[
(2+3)\cdot4
\]
Según la interpretación del producto, estamos sumando \(4\) un total de \(2+3\) veces:
\[
\underbrace{4+4}_{2\text{ veces}}
+
\underbrace{4+4+4}_{3\text{ veces}}
\]
Podemos reconocer entonces
\[
2\cdot4+3\cdot4
\]
Por lo tanto,
\[
(a+b)\cdot c
=
a\cdot c+b\cdot c
\]
Ambas formas expresan la misma idea:
\begin{gather*}
a\cdot(b+c)=a\cdot b+a\cdot c \\
(a+b)\cdot c=a\cdot c+b\cdot c
\end{gather*}
La propiedad distributiva será especialmente importante porque permite transformar expresiones que contienen productos y sumas sin modificar el número que representan.

\subsection{¿Qué hemos obtenido a partir de la suma?}

Resulta útil observar lo que acabamos de hacer.

Partimos de una interpretación muy sencilla:
\[
a\cdot b
=
\underbrace{b+b+\cdots+b}_{a\text{ veces}}
\]
A partir de ella pudimos comprender varias propiedades del producto:
\begin{gather*}
  a\cdot b=b\cdot a,\\
  (a\cdot b)\cdot c=a\cdot(b\cdot c),\\
  a\cdot1=a,\\
  a\cdot0=0,
\end{gather*}
y
\[
a\cdot(b+c)=a\cdot b+a\cdot c.
\]
Esto muestra que la relación entre suma y producto es más profunda que una simple abreviatura de escritura. Sobre los números naturales, el producto puede construirse a partir de sumas repetidas y gran parte de su comportamiento puede comprenderse utilizando las propiedades que ya conocemos de la adición.

Sin embargo, suma y producto continúan siendo operaciones diferentes. En una expresión es importante reconocer cuál de ellas está actuando en cada nivel.

Por ejemplo, en
\[
2+3\cdot4+5
\]
los sumandos de la expresión principal son
\[
2,\qquad3\cdot4,\qquad5,
\]
mientras que \(3\) y \(4\) son los factores del producto
\[
3\cdot4
\]
Esta distinción será fundamental para interpretar correctamente expresiones que combinen varias operaciones.
